\chapter{Conclusions and Future Work}
\label{chap:conclusions}

This thesis presented, to the best of our knowledge, the first approximate
nearest-neighbor search implementation for Tenstorrent hardware: an IVF-Flat
pipeline on a Wormhole n150 that combines device-side coarse search, host-side
scheduling of the lists selected by each 32-query batch, and distributed fine
search with hierarchical top-\(k\) selection. It was evaluated on
glove-100-angular against Faiss on two CPU configurations and an NVIDIA A100
system, at matched Recall@10, in terms of throughput, workload energy, and the
cost of each stage of the search path. This chapter answers the research
questions, summarizes the limitations, and outlines future work.

\section{Answers to the Research Questions}
\label{sec:rq-answers}

\paragraph{RQ1: Which ANN index is most suitable for Tenstorrent?}
IVF-Flat. Its fine search reduces to dense multiplications between query tiles
and contiguous blocks of database vectors, which map directly onto Tensix tile
operations and explicit DRAM transfers, whereas graph-based indexes depend on
pointer chasing and query-dependent control flow
(Chapter~\ref{chap:background}). This is an argued choice rather than a
measured comparison, since no other index was implemented. The implementation
showed that the regular part of the search maps well onto the device, but also that the irregular parts of
IVF, namely list selection, list sizes, and scheduling, remain the difficult
part of the mapping: they were placed on the host and they surface as worker
imbalance on the device.

\paragraph{RQ2: Can IVF search be implemented efficiently on Tenstorrent?}
It can be implemented, but not yet efficiently. The pipeline reaches 95.4\%
Recall@10 and runs at 5,550~queries per second at that recall, and grouping
queries into 32-query batches lets the fine search use full tiles. The same
batching, however, makes every query in a batch scan the union of the lists
selected by the whole batch. Because consecutive queries share almost no lists,
each query is scored against 12 to 32 times as many vectors as its own lists
contain, and against 81.7\% of the database at the 95\% operating point. This
raises recall by up to 7.5 percentage points at equal parameters, but at high
recall it turns IVF search into a nearly exhaustive scan. The union itself
saves data movement, since each distinct list is fetched once per batch
instead of once per query that requests it, but with queries in dataset order
this saving is small. Inside the
fine search, one page of candidates costs about 7~\(\mu\)s of compute, which is
far more than its four tile multiplications, and in the profiled window compute
waits for DRAM fetches for about 40\% of the time because the reader runs only
about one page ahead.

\paragraph{RQ3: What recall and throughput does Tenstorrent achieve?}
The A100 is the fastest system at every recall level, by 31.9 times at 60\%
recall and 7.8 times at 95\%. Tenstorrent is slower than both CPU
configurations up to about 88\% recall; at the 60\% target this is largely
due to a fixed cost of about 0.58~ms per batch that caps its throughput at low
\(N_{\mathrm{probe}}\). Its throughput, however, declines more slowly with
search depth: at 95\% recall it is 1.74 times faster than the 56-core CPU and
on par with the 112-core CPU. These CPU comparisons hold against the evaluated
Faiss grid, whose only point near 95\% recall uses 512 lists; Faiss
configurations with more lists, which were not evaluated, could narrow or reverse them.

\paragraph{RQ4: How does energy compare, and what are the limits of the
measurement?}
For a complete 100,000-query workload including initialization, the A100
system uses the least energy at every recall target, 9.3--20.1~mJ per query.
The Tenstorrent system uses 30.5--46.9~mJ per query, 2.3--3.8 times the A100
value, but 53--79\% less than the CPU configurations, which use
69--219~mJ. The Wormhole device itself draws only 24--34~W and accounts for
11--18\% of the Tenstorrent system energy; the rest is drawn by its dual-socket
host. The main limits of the measurement are that it includes initialization,
so per-query values depend on workload length; that the device sensors differ,
chip telemetry on Tenstorrent and board power on the A100; and that the three
systems use different hosts.

\paragraph{RQ5: Which stages dominate, and why?}
Fine search accounts for 90--98\% of the Tenstorrent search time at
\(N_{\mathrm{list}}=512\), the only list count for which the stages were
timed; host preparation, transfers, coarse search, and
output handling together take only 26--30~ms per 10,000 queries. Within fine
search, batch transitions and the final aggregation overlap with the workers
and are not on the critical path at these configurations; at low
\(N_{\mathrm{probe}}\), however, the serial aggregation of 63 worker results
per batch is a likely cause of the fixed per-batch floor. The time is instead spent in the per-list
loop of each worker. With at most 21 floating-point operations per fetched
byte, and only 0.7--1.7 useful ones, the search is memory-bound in the roofline
sense and uses only 2\% of the peak arithmetic rate; its time is set by
irregular DRAM fetch times with shallow buffering, which create pipeline
bubbles, and by the work done on every page around the multiplication, most
likely the top-\(k\) update. Batch completion is
further set by the slowest worker, because tasks remain coarse: whole lists,
or chunks of up to 128 pages for longer lists, which occur only with 512
lists. In a schedule without chunking, which matches the timed runs at 1024
and 2048 lists, the busiest worker holds a median of 1.3--4.2 times the mean
load, the most at low \(N_{\mathrm{probe}}\), where fewer lists than workers
are available. More than any of these effects, however, the batch union
multiplies the number of pages that the workers must process.

\paragraph{The energy hypothesis.}
The motivation of this thesis was that an energy-efficient machine-learning
accelerator might also provide lower-energy vector search. At the system level,
this is not confirmed for the evaluated setup: the A100 system uses less energy
at every recall target. The measurements do, however, locate the cause. The
Wormhole device is frugal: at the 95\% target the chip uses 8.5~mJ per query,
similar to the 10.7~mJ of the A100 board, while the system energy is dominated
by the host and by the time spent in the search loop. Both can be reduced,
by a lighter host and by shorter sessions, which motivates the future work
below.

\section{Limitations}
\label{sec:conclusion-limitations}

The implementation is limited by the platform and by its own design in four
ways.

\paragraph{Limited low-level primitives.} The low-level kernel (LLK) API of
TT-Metalium offers little room to customize the operations that the search
relies on. Ordering scores together with 32-bit identifiers is possible only
through a specific top-\(k\) path, which in turn requires FP32 destination
accumulation (Section~\ref{subsec:tiled-list-layout}). Other primitives, such
as the merge of sorted results, support only 16-bit indices, which cannot
identify the 1.18 million database vectors directly. The design therefore
follows what the available primitives allow rather than what the algorithm
would ideally use.

\paragraph{Device memory.} The complete index must reside in the device DRAM
of the Wormhole n150, 12~GB \cite{tenstorrent_wormhole_spec}. With vectors
padded to 128 FP16 dimensions and one \(32\times32\) identifier tile per page,
each database vector occupies 384~bytes, which bounds a single device to
roughly 30 million vectors, far from billion-scale collections. Larger collections could in
principle be swapped in and out of the device, as tiling does between DRAM and
L1, but the lists would then cross the PCIe link, about 32~GB/s against
288~GB/s for the device DRAM. Since a batch union covers a large part of the
database, about 374~MB at \((512,32)\), swapping pays off only when batches
touch few lists, for example after grouping similar queries, or when the most
frequently used lists are kept on the device. Such out-of-core search is
usually treated as a separate problem from in-memory search: systems such as
DiskANN \cite{diskann} and SPANN \cite{spann} keep most of a billion-scale
index on SSDs and are designed around that storage, and the NeurIPS'21
billion-scale challenge evaluated in-memory and out-of-core solutions in
separate tracks \cite{bigann_neurips21}. This thesis addresses the in-memory
setting.

Within the device, the L1 memory
of each Tensix core, shared by the circular buffers, query tiles, and top-32
states, in turn limits how deeply the reader can buffer ahead of compute.

\paragraph{Load imbalance.} Tasks are coarse: lists are split only when they
exceed 128 pages, so a worker that receives a large task can still set the
completion time of the whole batch, and at low \(N_{\mathrm{probe}}\) fewer
tasks than workers are available.

\paragraph{Partial last pages.} The last page of an inverted list is usually
only partially filled. Before the top-\(k\) update, its padded positions must
receive a very low similarity, which the compute kernel achieves by scanning
all \(32\times32=1{,}024\) entries of the page
(Section~\ref{subsec:tiled-list-layout}). Every list in a batch union pays
this cost once, so it weighs most when the lists are small, as with 2048
lists.

The platforms differ in host processors and hardware generations,
so the comparison is between complete systems rather than processor
architectures. Energy includes initialization and uses different device
sensors on the two accelerators.

\section{Future Work}
\label{sec:future-work}

The results point to one main direction: an index designed specifically for
Tenstorrent, which is expected to raise its throughput more than refinements
of the current design, because the IVF scan fetches far more data than each
query needs and is limited by per-page work and irregular fetch latency rather
than by arithmetic. Within the current IVF design, several further
steps would help: fetching less data per batch, distributing it better over
the workers, and making each page cheaper to process.

\begin{enumerate}
    \item \textbf{Design an index for the device.} The main direction, and
    the one that could raise throughput on the device the most, is an index
    designed specifically for the hardware, as CAGRA did for GPUs
    \cite{cagra}. A graph-based index can reach high recall
    while visiting far fewer vectors than an IVF scan, but, as discussed in
    Section~\ref{sec:index-suitability}, a direct port of a CPU graph
    algorithm would not suit Tenstorrent. A graph designed for the device
    would take its constraints into account from the start: the limited L1
    memory of each core, the tile granularity of computation, and the cost of
    DRAM fetches shared by many cores.
    \item \textbf{Group queries before batching.} In dataset order the queries
    of a batch share almost no lists. Sorting them by their coarse selections
    before batching makes each batch fetch fewer distinct lists: in a first
    test at \(N_{\mathrm{list}}=512\), with 8 partitions, fine search drops by
    57\% at \(N_{\mathrm{probe}}=16\) and by 42\% at \(N_{\mathrm{probe}}=32\),
    for a host sorting cost of 100--180~ms per 10,000 queries. Its effect on
    recall remains to be measured.
    \item \textbf{Use only lightweight quantization.} Quantization reduces the
    bytes fetched per vector, but only simple formats suit the device.
    Product quantization and other codebook-based schemes need a lookup table
    to convert each code back to a value, which does not map well onto tile
    operations. Scalar formats that can be converted directly, such as 8-bit
    or 4-bit values, are the practical options.
    \item \textbf{Relieve the DRAM contention.} A high fetch volume is costly
    on this device because the 63 workers share six DRAM controllers with
    twelve GDDR6 channels \cite{tenstorrent_wormhole_spec}, so their readers
    compete for the same channels, which is consistent with the slow fetches
    observed in the single-list trace. Deeper buffering does not help: a
    first test with an eight-page instead of a two-page reader prefetch made
    fine search about 7\% slower, as the additional outstanding reads add to
    the contention. Reducing the fetched volume is therefore the main remedy,
    and repeating the page-level trace at larger \(N_{\mathrm{probe}}\)
    would quantify the contention directly.
    \item \textbf{Distribute the work in finer slices.} The implementation
    already splits lists longer than 128 pages into chunks. Finer slices, down
    to a few pages, would balance the load further across the workers, so that
    no single large task sets the completion time of a batch. The single
    aggregation core would then remain a bottleneck only at low
    \(N_{\mathrm{probe}}\). A tree-shaped reduction is not expected to be an efficient
    alternative: intermediate reduction nodes would need L1 space that the
    fine search uses for its fetch buffers, and every core would have to act
    both as a worker and as a reduction node.
    \item \textbf{Scale beyond one chip and one dataset.} Datasets larger than
    device memory, other embedding models, and multi-chip Tenstorrent systems
    would test whether the approach holds at the scale of heavy retrieval
    workloads.
    \item \textbf{Integrate retrieval with inference.} The original motivation,
    running retrieval on the same hardware as the language model in a RAG
    pipeline, remains open. Co-locating both stages on one device would remove
    a transfer between systems and is the setting in which a Tenstorrent
    retrieval stage would be most valuable.
\end{enumerate}

\cleardoublepage
